\chapter{十分之几与百分之几}\label{ch:g4:decimals}

在 $3$ 与 $4$ 之间，整数便“沉默”了——
但一位跑者可能已领先 $3$ 米多一点。把单位平均分成十份，再分成一百份，
便得到了\emph{小数}，以及表示它们的一种巧妙写法——
\omterm{def:g4:decimals:point}{小数点}。（小数的完整内容
在\cref{ch:g5:decimals}中。）

\section{小数}

\begin{definition}[十分之一与百分之一]\label{def:g4:decimals:def}
把一个单位平均分成 $10$ 份，得到\emph{十分之一}，
$\frac{1}{10}$；再把每个十分之一分成 $10$ 份，得到
\emph{百分之一}，$\frac{1}{100}$。于是
\[
1 = \frac{10}{10} = \frac{100}{100},
\qquad
\frac{1}{10} = \frac{10}{100}.
\]
\end{definition}

\begin{omfigure}
\begin{tikzpicture}[scale=0.42]
  \draw[gray!50, thin] (0,0) grid (10,10);
  \foreach \i in {0,...,9} \fill[omDef!35] (3,\i+0.06)
    rectangle (3.94,\i+0.94);
  \fill[omThm!70] (6.05,4.05) rectangle (6.95,4.95);
  \draw[thick] (0,0) rectangle (10,10);
  \node[below, font=\small] at (5,-0.4)
    {整个正方形 $= 1$；\ \ 一列（蓝色）$= \frac{1}{10}$；\ \
    一个小格（红色）$= \frac{1}{100}$};
\end{tikzpicture}

{\small 百格图：共十列，每列十格。一列是
十分之一，一格是百分之一。}
\end{omfigure}

\begin{example}\label{ex:g4:decimals:count}
在百格图中涂色 $3$ 列和另外的 $7$ 个小格，涂色部分为
\[
\frac{3}{10} + \frac{7}{100} = \frac{30}{100} + \frac{7}{100}
= \frac{37}{100}
\]
即整个正方形的——三十七个百分之一。
\end{example}

\section{小数点}

\begin{definition}[小数的写法]\label{def:g4:decimals:point}
数 $2 + \frac{3}{10} + \frac{7}{100}$ 可以写成
\[
2.37
\]
——\emph{小数点}\index{小数点}把整数部分（$2$）与
小数部分分开；小数点后第一位数字表示有几个十分之一，
第二位数字表示有几个百分之一。读作“二点三七”，或“二又
三十七个百分之一”。
\end{definition}

\begin{example}\label{ex:g4:decimals:convert}
\[
\frac{7}{10} = 0.7, \qquad
\frac{43}{100} = 0.43, \qquad
\frac{5}{100} = 0.05, \qquad
3 + \frac{4}{10} = 3.4 .
\]
注意 $\frac{5}{100}$：五个\emph{百分之一}需要在
十分位上写 $0$——写作 $0.05$，而不是 $0.5$。
\end{example}

\begin{omfigure}
\begin{tikzpicture}[scale=1.35]
  \draw[-Stealth] (-0.2,0) -- (8.4,0);
  \foreach \i in {0,10,20,30,40,50,60,70,80}
    \draw ({\i*0.1},-0.09) -- ({\i*0.1},0.09);
  \foreach \i in {0,...,80} \draw ({\i*0.1},-0.045) -- ({\i*0.1},0.045);
  \foreach \x/\l in {0/2, 1/{2.1}, 2/{2.2}, 3/{2.3}, 4/{2.4},
                     5/{2.5}, 6/{2.6}, 7/{2.7}, 8/{2.8}}
    \node[below, font=\footnotesize] at (\x,-0.1) {$\l$};
  \fill[omThm] (3.7,0) circle (1.7pt)
    node[above, font=\small] {$2.37$};
\end{tikzpicture}

{\small 放大观察 $2$ 与 $2.8$ 之间：大刻度表示每个十分之一，
小刻度表示每个百分之一。数 $2.37$ 位于 $2.3$ 与
$2.4$ 之间的第 $7$ 个小刻度处。}
\end{omfigure}

\begin{example}[钱与百分之一]\label{ex:g4:decimals:money}
一分钱是一欧元（或一美元、一英镑）的百分之一：价格
为 $4.85$ 表示 $4$ 枚完整的硬币和 $85$ 个百分之一。你每天用的钱，
就用到了\omterm{def:g4:decimals:point}{小数点}：$10$ 分钱 $= 0.10$ 元；两欧元五十分
$= 2.50$ 欧元；而 $0.05$ 是一枚五分硬币，不是五十！
\end{example}

\section{用十分之几比较大小}

\begin{method}[比较小数（初步认识）]\label{met:g4:decimals:compare}
\begin{enumerate}
  \item 先比较整数部分：$3.2 > 2.9$。
  \item 整数部分相同时，先比较十分位，再比较百分位：
        $2.37 < 2.5$，因为 $3$ 个十分之一 $< 5$ 个十分之一。
  \item 把两个数都写成两位小数会更清楚：$2.5 = 2.50$，
        即 $37 < 50$ 个百分之一。
\end{enumerate}
易错之处：$2.37$ 比 $2.5$ 位数多，但
\emph{更小}——位数多并不代表数更大！
\end{method}

\begin{example}\label{ex:g4:decimals:compare}
将 $0.4$、$0.35$ 和 $0.09$ 排序：统一写成两位小数，得 $0.40$、$0.35$、
$0.09$；因此
\[
0.09 < 0.35 < 0.4 .
\]
\end{example}

\section{练习}

\begin{exercise}[$\star$]\label{exo:g4:decimals:1}
在百格图中，若涂色的是：$5$ 列；
$2$ 列和 $3$ 格；$45$ 格。涂色部分各占多少？用\omterm{def:g4:fractions:def}{分数}作答。
\end{exercise}

\begin{exercise}[$\star$]\label{exo:g4:decimals:2}
用小数表示：
$\frac{9}{10}$; \quad $\frac{27}{100}$; \quad $\frac{3}{100}$; \quad
$5 + \frac{6}{10}$; \quad $\frac{60}{100}$.
\end{exercise}

\begin{exercise}[$\star$]\label{exo:g4:decimals:3}
写成\omterm{def:g4:fractions:def}{分数}（十分之几或百分之几）：
$0.3$; \quad $0.51$; \quad $0.07$; \quad $1.9$.
\end{exercise}

\begin{exercise}[$\star$]\label{exo:g4:decimals:4}
在 $6.58$ 中，数字 $5$ 表示的是多少个什么？数字 $8$ 呢？
仿照\cref{def:g4:decimals:point}，将该数分解。
\end{exercise}

\begin{exercise}[$\star$]\label{exo:g4:decimals:5}
在从 $4$ 到 $5$、以十分之一为刻度的数轴上标出：
$4.2$；\ $4.75$（位于哪两个刻度之间？）；\ $4.9$。
\end{exercise}

\begin{exercise}[$\star$]\label{exo:g4:decimals:6}
抄写并在空格处填入 $<$、$>$ 或 $=$：
\[
0.6 \;?\; 0.60, \qquad
2.8 \;?\; 2.75, \qquad
0.09 \;?\; 0.1, \qquad
5.3 \;?\; 5.13 .
\]
\end{exercise}

\begin{exercise}[$\star$]\label{exo:g4:decimals:7}
从小到大排序：$1.05$；\ $1.5$；\ $0.95$；\ $1.45$。
\end{exercise}

\begin{exercise}[$\star$]\label{exo:g4:decimals:8}
用小数写出：三欧元五分；十二欧元
五十分；九十九分。然后将这三个价格排序。
\end{exercise}

\begin{exercise}[$\star$]\label{exo:g4:decimals:9}
哪些一位小数严格位于
$3$ 与 $4$ 之间？请全部列出。共有多少个？
\end{exercise}

\begin{exercise}[$\star\star$]\label{exo:g4:decimals:10}
汤姆说：“$0.25$ 比 $0.5$ 大，因为 $25$ 比
$5$ 大。”请借助百格图向他指出错误（
每个数各涂了多少格？）。
\end{exercise}

\begin{exercise}[$\star\star$]\label{exo:g4:decimals:11}
一名运动员三次跳远的成绩分别是 $3.6$ m、$3.58$ m 和 $3.65$ m。
将这三次成绩从短到长排序。哪两次相差
恰好为 $\frac{5}{100}$ 米？
\end{exercise}
